
The study of submanifolds embedded in a manifold is a classical subject
in geometry, analysis and topology, cf. \cite{dajczer2019submanifold},\cite{lee2012smooth},\cite{tu2017differential}. The general
framework
has both intrinsic aspects, e.g. the study of geodesics, connections and
curvature
of embedded submanifolds, as well as the study of extrinsic properties
encapsulated by the
shape of embedded submanifold including the second fundamental form, mean
curvature, normal bundle, etc.

Denoting by $\iota: M\hookrightarrow X$ the smooth embedding of a manifold
$M$ into ambient manifold $X$, there is a fundamental short exact sequence
$$
0 \longrightarrow \tangent{M} \longrightarrow \iota^*\tangent{X} \longrightarrow \normal{\iota}
\longrightarrow 0
$$
of vector bundles on $M$, where $\normal{\iota}$ denotes the normal vector bundle of
$\iota$.
This allows to introduce the notion of tubular neighborhood and exponential
map, whose domain is the open disk bundle in $\normal{\iota}$. On the topological side of
the problem, the tangent and the normal bundles of $M$ interact with each other 
through the ambient manifold $X$, e.g. when $X$ is a Euclidean space then 
$\tangent{M}$ and $\normal{\iota}$ stably equivalent with implications on their characteristic 
classes. As an exposition of the theory of characteristic classes, we recommend \cite{hatcher2003vector},\cite{milnor1974characteristic}, and, for an approach more along the lines of differential geometry, we refer to \cite{tu2017differential}.

The motivating example is the Euler characteristic class of the normal
bundle
of embedded oriented submanifold $M$ of $X$. This class has values in the
cohomology group $\mathrm{H}^{\mathrm{codim}(M)}(M,{\mathbb Z})$, with $\mathrm{codim}(M)$ denoting
the codimension of $M$ in $X$. It captures several properties of
the pair $M, X$: it measures the obstruction to finding a nowhere vanishing
normal vector field, and also governs the local intersection behavior of the
embedding via the Thom class of the normal bundle over $M$.

The main focus of our article is the computation of characteristic
classes
of normal vector bundles for certain class of pairs $M, X$. In particular, the
ambient
manifold $X$ is a homogeneous space called generalized flag manifold, and the
submanifold $M$ is an embedded homogeneous submanifold of the same geometrical
type.
In practice, this means that for a complex simple Lie group $G_n$ belonging to
one of the classical series $A_n, B_n, C_n, D_n$ of simple Lie groups (see \cite{fulton2013representation},\cite{goodman2000representations}), and
$P_n\subset G_n$ its standard parabolic subgroup, $X=G_n/P_n$ and $M=G_{n-1}/P_{n-1}$
together with the compatibility condition $P_{n-1}=G_{n-1}\cap P_n$ satisfied.
The tangent bundles of the
generalized flag manifold $G_n/P_n$ and its embedded submanifold
$\iota : G_{n-1}/P_{n-1}\hookrightarrow G_n/P_n$
fit into a short exact sequence
$$
0 \longrightarrow \tangent{(G_{n-1}/P_{n-1})} \longrightarrow \iota^*\tangent{(G_n/P_n)}
\longrightarrow \normal{\iota} \longrightarrow 0
$$
of vector bundles on $G_{n-1}/P_{n-1}$ with $\normal{\iota}$ denoting the normal bundle of $\iota$.
An illustrative example consists of the embedded $\cp{n-1}$ in $\cp{n}$ induced 
by the standard block-diagonal embedding of $\gl{n}$ in $\gl{n+1}$, with the 
normal bundle being the dual of the canonical line bundle over $\cp{n-1}$,
whose first Chern class generates the cohomology ring of $\cp{n-1}$.

It is well known that the category of algebraic or holomorphic vector
bundles on
a generalized flag manifold $G/P$ is equivalent to the category of
finite-dimensional
linear representations of the parabolic subgroup $P$. Consequently, every
algebraic or
holomorphic vector bundle on a generalized flag manifold is equivalent to
a homogeneous
vector bundle. In addition, geometrical problems for generalized flag
manifold can be
reformulated in terms of representation theory of the pair $G$ and $P$, hence
the characteristic classes of homogeneous vector bundles are controlled by
the
representation theory associated to symmetries related to the generalized
flag
manifold. The general concept of homogeneous vector bundles from the point 
of view of representation theory is studied in the seminal papers 
\cite{atiyah1961vector},\cite{borel1958characteristic},\cite{bott1957homogeneous}.
In fact, our paper can be regarded as a substantial step towards completion 
of their program by giving explicit description of Chern classes 
of normal bundles of embeddings of the form $\iota : G_{n-1}/P_{n-1}\hookrightarrow G_n/P_n$.

The algebraic nature of classical simple Lie groups allows one to exploit 
the tool of algebraic geometry known as the Schubert calculus - a geometrical 
approach to enumerative geometry, which develops combinatorics and 
intersection theory of Schubert cells in Grassmannians.
Schubert calculus is concerned with the ring structure on the cohomology
of flag
varieties and Schubert varieties and has been explored in depth in, e.g., 
\cite{bernstein1973schubert},\cite{hiller1986pieri},\cite{pragacz1994pieri},\cite{pragacz2003pieri}. (For a briefer and more compact exposition, we refer to \cite{naughton2021schubert}.)
% A key result is that the Schubert cells,
%or rather,
%the classes of their Zariski closures - the Schubert cycles or Schubert
%varieties - span
%the whole cohomology ring. 
%Several authors have since generalized the techniques of Schubert calculus to quantum cohomology of complex flag manifolds in, for example, %\cite{buch2003gromov},\cite{kresch2003quantum},\cite{kresch2004quantum}. 
In this perspective, our results can be regarded as a relative version of Schubert calculus 
for pairs of generalized flag manifolds $G_{n}/P_{n}$ and $G_{n-1}/P_{n-1}$.

We remind the reader that the present article only treats the case of pairs $(G_{n}/P_{n}, G_{n-1}/P_{n-1})$ with $G_{n}$ and $G_{n-1}$ of the same Lie type. Moreover, both groups $G_{n}, G_{n-1}$ are complex. Thus, there are several ways in which we can expand on the present article in the future. For one, we may consider pairs $(G_{n}/P_{n}, G_{n-1}/P_{n-1})$ with $G_{n}$ and $G_{n-1}$ of different Lie types, and perhaps even the exceptional types $E_{6}, E_{7}, E_{8}$ or $F_{4}$ or $G_{2}$. Another possibility is to consider real and quaternionic flag manifolds $X$ and to compute the $\coh{*}{X; \mathbb{Z}_{2}}$-valued Stiefel-Whitney classes or $\coh{*}{X; \mathbb{Z}}$-valued Pontryagin classes of normal bundles of their embeddings. 

As for the applications of our results, we already mentioned the relative Schubert calculus. Another arena
where our results should be useful consists of various questions and problems in differential and algebraic geometry
of flag manifolds and their homogeneous submanifolds, e.g. the vector-valued problems for equivariant operators
with the initial data on homogeneous submanifolds. 
%The next natural step is to consider analogous questions 
%and problems for a general pair of a curved manifold and its submanifold whose Klein models are %studied in the 
%present article.

Our article is organized as follows. In order to have a self-contained introduction that allows for a
smooth presentation of our original results in later sections, Section 2 starts with a review of the
classical series of complex simple Lie groups $A_{n}, B_{n}, C_{n}$ and $D_{n}$. Then we discuss 
homogeneous complex vector bundles over homogeneous spaces $G_{n}/P_n$, with particular emphasis 
on their representation-theoretical aspects as well as their splitting properties when restricted 
to an embedded submanifold $G_{n-1}/P_{n-1} \hookrightarrow G_{n}/P_{n}$. We use the representation-theoretical tools to describe normal vector bundles for these embeddings induced by the standard 
inclusion of 
$G_{n-1}$ in $G_{n}$. The bundles split into direct sums of so called canonical 
and complementary canonical bundles over $G_{n-1}/P_{n-1}$, which are pull-backs of certain
homogeneous vector bundles over linear and isotropic Grassmanians.
The section concludes with a description of the realization of the cohomology ring of $G_{n-1}/P_{n-1}$ 
as a tensor product of cohomology rings of some particular Grassmannians. This then reduces the initial task 
of calculating the Chern classes of normal vector bundles over $G_{n-1}/P_{n-1}$ to 
the case when $G_{n-1}/P_{n-1}$ themselves are (linear or isotropic) Grassmannians. 

In Section 3, we briefly recall the well-known CW structure and cohomology of linear complex 
Grassmannians $\gr{n}{k}$, whose cells and cohomology generators are in bijective correspondence 
with a collection of certain partitions. We study the impact of various embeddings and diffeomorphisms 
between linear Grassmannians, which will also find their use in Section 4, on the cohomology generators 
of $\coh{*}{\gr{n}{k}}$. This in turn enables an
inductive calculation of Chern classes of normal bundles of embeddings of linear Grassmannians.

In Section 4, we obtain a CW structure on the isotropic Grassmannians (series $B_{n}, C_{n}, D_{n}$) by inspecting their embeddings in linear Grassmannians. The knowledge of the CW structure on the isotropic 
Grassmannians immediately yields the generators of their cohomology rings corresponding to partitions 
with some additional properties imposed by the isotropy conditions. Afterwards, we discuss the effects 
of specific embeddings and other particular maps between the isotropic Grassmannians on these generators. These observations 
are then used for calculations of Chern classes of normal bundles of embeddings of isotropic 
Grassmannians.